\scanpage{26}
\mspage{20}

Notons maintenant que b) n'est autre que
la variante affaiblie de a), quand on omet
de spécifier que $C_0,TC_0$ doivent être \unsure{petites|petits}.
Donc le seul problème pour $b\Rightarrow a$,
c'est de voir ceci :

\uline{Lemme conjectural} \textcircled{1} \deleted{\illegible{2}}
\deleted{Si on a une partie
close à gauche $\unsure{\Phi|\underline{\varphi}|\underline{\mathcal{Q}}|\underline{\mathcal{D}}}$ de $\mathrm{Fl}\,M$, telle que
$\Phi_*\subset\mathrm{Fib}_W$ (resp. $\Phi_*\subset W$), alors il
existe une partie \uline{\unsure{petite}} $\Phi_0\subset\Phi$,
telle que $\Phi_{0*}\subset\mathrm{Fib}_W$ (resp. $\Phi_{\unsure{0|0*}}\subset W$).}
% slanted
\margin{Ce lemme a l'air un peu faux, mais ça doit pouvoir s'arranger\ldots}

\inserted{\textcircled{2} Si $C,C'$ sont deux parties closes à gauche de $\mathrm{Fl}\,M$ et
$C$ et $C'$ sont accessibles, il en est de \unsure{$=$} de $C\cap C'$, \unsure{Si $a''$)|il y a}}

Je crois que c'est \deleted{\unsure{bel}}
sous les hypothèses faites sur $M,W$ ---
qui doivent impliquer que dans un sens
convenable, $W$ est une partie ``accessible''
de l'ens $\mathrm{Fl}\,M$ des flèches de la cat.
accessible $M$. Il doit en résulter qu'une
partie de $M$ construite explicitement en
terme de $W$, telle $\mathrm{Fib}_W$, est elle-$=$
accessible, et le lemme doit sortir de
là.

Ainsi, $2^\circ$ est prouvé, modulo ce lemme
d'accessibilité, qu'il va falloir prouver
$\pm$ en $=$ \deleted{temps} que le th.~1 (que je n'ai
fait \unsure{ici} qu'énoncer) ;
% archive sheet number: 25
